\chapter{双対空間}
    \section{基底変換と座標変換の関係}
        \begin{shadebox}
            $V$を$n$次元ベクトル空間とし($\realNumbers^n$や$\complexNumbers^n$とは限らないことに注意)，$\{\vv_1,\dotsc,\vv_n\}$と$\{{\vv_1}',\dotsc,{\vv_n}'\}$を$V$の基底とする。
            これらが$n$次正則行列$A$により次の関係式に従うとする。
            \[
                \begin{bmatrix}
                    {\vv_1}' \\
                    \vdots \\
                    {\vv_n}'
                \end{bmatrix}
                = A
                \begin{bmatrix}
                    \vv_1 \\
                    \vdots \\
                    \vv_n
                \end{bmatrix}
            \]
            $\{\vv_1,\dotsc,\vv_n\},\{{\vv_1}',\dotsc,{\vv_n}'\}$の双対基底をそれぞれ$\{\vv^1,\dotsc,\vv^n\}\,\{{\vv^1}',\dotsc,{\vv^n}'\}$とし，
            $\underline{\vx} \coloneq [x_1,\dotsc,x_n]^\top,\underline{\vx}' \coloneq [{x_1}',\dotsc,{x_n}']^\top$とする。
            このとき$[{\vv^1}',\dotsc,{\vv^n}']\underline{\vx}' = [\vv^1,\dotsc,\vv^n]\underline{\vx}$であれば$\underline{\vx}' = A\underline{\vx}$である。
            すなわち，双対空間での座標変換と元の空間の基底間の線形変換が同じ形になる。
        \end{shadebox}
        \begin{proof}
            $[{\vv^1}',\dotsc,{\vv^n}']\underline{\vx}' = [\vv^1,\dotsc,\vv^n]\underline{\vx}$より
            \begin{align*}
                \underline{\vx}' &= 
                \begin{bmatrix}
                    {\vv_1}' \\
                    \vdots \\
                    {\vv_n}'
                \end{bmatrix}
                [\vv^1,\dotsc,\vv^n]\underline{\vx} = A
                \begin{bmatrix}
                    \vv_1 \\
                    \vdots \\
                    \vv_n
                \end{bmatrix}
                [\vv^1,\dotsc,\vv^n]\underline{\vx} = A I_3\underline{\vx} = A\underline{\vx}
            \end{align*}
        \end{proof}